

In the spirit of section \ref{section: yoga} we will work towards constructing a $6$-functor formalism suitable to imply coherent duality. Contrary to the classical approach we will not associate to a scheme $X$ the derived category of (quasi-)coherent $\struct{X}$-modules but instead enlarge this category considerably. We will -- in case of an affine $X=\Spec A$ -- associate to $X$ the derived category $\derived{A_\blacksquare}$ of modules \emph{complete} over the \emph{analytic ring} $A_\blacksquare$. Defining the notion of (pre-)analytic rings the associated notion of completeness is the content of this chapter. The specific (pre-)analytic rings $A_\blacksquare$ will be introduced in section \ref{section: the pre-analytic rings Ablack and (A,R)black}, their analyticity will be shown in sections \ref{section: analyticity of Ablack and (A, R)black} and \ref{section: an analogue of hilbert's basis theorem}.

\section{Pre-analytic \& analytic rings}
\label{section: pre-analytic and analytic rings}

What is a (pre)-analytic ring supposed to be? During their investigation of condensed mathematics, Clausen and Scholze noticed that there is a reoccurring theme: Inside a category of condensed modules one finds a certain subclass of \quote{nice} objects one deems \quote{complete}. This covers for example the case of so-called \emph{solid} modules (the ones we will study in this thesis) but also various categories of analytically interesting modules like $p$-liquid vector spaces inside condensed $\R$-modules (which we will neither define nor investigate). A common theme is that one has the following data: Certain \emph{free complete} modules $\pfree{A}{S}$ associated to any extremally (or even profinite) set $S$ and for any such $S$ a natural morphism $S\to \pfree{A}{S}$ of (condensed) sets. This is what one call a \emph{pre-analytic ring}. Any module $M$ in question is then said to be \emph{complete} if for any $S$ and any specification of how elements of $S$ map into $M$ (\ie a map $S\to M$ of (condensed) sets) there is a unique extension of the map to $\pfree{A}{S}$. Depending on the free complete modules chosen this provides varying flavors of theory. In the solid case the theory is very non-archimedean, however in the $p$-liquid case the theory is well suited to \quote{traditional} functional analysis over the real numbers. As the notion of pre-analytic rings is very general, one has to impose some regularity to obtain a good theory of completeness (for example that the full subcategory of complete modules being abelian). This then leads to the notion of \emph{analyticity}, which for example implies that the free complete modules $\pfree{A}{S}$ are themselves complete (something that is not a priori clear). For an \emph{analytic ring} (a pre-analytic ring that is analytic) a myriad of good results follow at once, \confer theorem \ref{theorem: properties of the category of complete modules}. This is the content of this section.

\begin{definition}[Pre-analytic rings]
	\label{definition: pre-analytic rings}
	
	A \emph{pre-analytic ring} $\preAnaRing A$ is a triple $(\underlying A, \pfree A -, \eta)$, where $\underlying A$ is a condensed ring, called the \emph{underlying ring}, $\pfree A - \colon \extremallyDisconnectedSets \to \mod{\underlying A}$ is a functor preserving finite disjoint unions, called the \emph{free completion} and $\eta$ is a natural transformation, containing the datum of a morphism $\eta_S\colon S \to \pfree A S$ of condensed sets for every $S\in \extremallyDisconnectedSets$.
	
	A \emph{map of pre-analytic rings} $\preAnaRing{A} \to \preAnaRing{B}$ is a pair $(\phi, \epsilon)$ where $\phi$ is a map of underlying rings $\underlying A \to \underlying B$ and $\epsilon$ is a natural transformation $\pfree A - \Rightarrow \scalarRestriction{\phi}\pfree B -$ consisting of maps of $\underlying A$-modules, such that for all extremally disconnected $S$ the induced diagram
	\begin{center}
		\begin{tikzcd}[column sep=0.75em]
								&S\ar[dl]\ar[dr]	&\\
			\pfree A S\ar[rr]	&					&\scalarRestriction{\phi}\pfree B S
		\end{tikzcd}
	\end{center}
	of condensed sets commutes.
	
	Denote the category of pre-analytic rings by $\preAnalyticRings$.
\end{definition}

\begin{example}
	\label{example: examples for pre-analytic rings}
	We have the following list of examples of pre-analytic rings:
	\begin{itemize}
		\item 
		Any condensed ring $A$ naturally admits the structure of a pre-analytic with underlying ring $A$, free completion $\free{A}{-}$ and natural map $S\to \free{A}{S}$ given by adjointness. This pre-analytic ring will be denoted by $A$ as well. We obtain a faithful embedding of condensed rings into $\preAnalyticRings$.
		
		\item
		The pre-analytic ring $\Z_\blacksquare$ with underlying ring $\Z$ and free completion $S=\limit_i S_i\mapsto \free{\Z_\blacksquare}{S}\ldef \limit_i \free{\Z}{S_i}$ along with the natural map $S\to \free{\Z_\blacksquare}{S}$ induced by the cone $S_i\to \free{\Z}{S_i}$.
		
		\item
		More general we will define for any discrete ring $A$ the pre-analytic ring $A_\blacksquare$ with underlying ring $A$ in section \ref{section: the pre-analytic rings Ablack and (A,R)black}. For any discrete $R$-algebra $A$ we will also define a relative version $(A, R)_\blacksquare$.
	\end{itemize}
\end{example}

\begin{definition}[Complete modules]
	\label{definition: completeness}
	
	Let $\preAnaRing{A}$ be a pre-analytic ring. An $\underlying{A}$-module $M$ is called \emph{$\preAnaRing{A}$-complete} if for any extremally disconnected set $S$ the natural map
		$$\Hom_{\underlying{A}}(\pfree A S, M) \to \Hom(\associated{S}, M),$$
	induced by $\associated{S}\to \pfree A S$, is an isomorphism. That is, for each solid diagram
	\begin{center}
		\begin{tikzcd}
			\associated{S} &[-1.5em]&[-1.5em] M\\
			&\pfree A S
			\ar[from=1-1, to=1-3, "\forall"]
			\ar[from=1-1, to=2-2]
			\ar[from=2-2, to=1-3, "\exists!\text{ $\underlying{A}$-linear}", swap, dotted]
		\end{tikzcd}
	\end{center}
	there is a dotted arrow making the triangle commute. Write $\pCompleteMod A$ for the (full) subcategory of $\mod{\underlying{A}}$ of $\preAnaRing{A}$-complete modules.
\end{definition}

\begin{remark}[Why \quote{completeness}?]
	This definition is of course rather abstract. We will see an interpretation of $\preAnaRing{A}$-completeness in the case of $\preAnaRing{A} = \Z_\blacksquare$ in remark \ref{remark: interpreting Zblack-completeness}. In this case the completeness of some (discrete) $\underlying{A}$-module $M$ requires \emph{certain sequences of elements in $M$ to \quote{converge}}. This then of course explains why this notion is termed \emph{completeness}.
\end{remark}

\begin{remark}
	Let $\preAnaRing{A}=(\underlying{A}, \pfree A -, \eta)$ be a pre-analytic ring. Since the functor $\ufree A -$ giving free $\underlying A$-modules is left adjoint to the forgetful functor $\mod{\underlying{A}} \to \condensedSets$, the natural map $\eta$ gives rise to morphisms $\ufree A S\to \pfree A S$ of $\underlying{A}$-modules, natural in the extremally disconnected set $S$.
	
	An $\underlying{A}$-module $M$ is then $\preAnaRing{A}$-complete if and only if the induced morphism
		$$\Hom_\underlying{A}(\pfree{A}{S}, M) \to \Hom_\underlying{A}(\ufree{A}{S}, M)$$
	is an isomorphism for all extremally disconnected $S$.
\end{remark}

We will now start to investigate when the category $\completeMod{\preAnaRing{A}}$ of $\preAnaRing{A}$-complete modules is well-behaved. Informally, if one is able to show that the free completions $\pfree{A}{S}$ are $\preAnaRing{A}$-complete (which is not clear!), then one can argue that all modules presentable by such free completes are $\preAnaRing{A}$-complete as well. 
\begin{definition}[Free and finitely free complete modules]
	Let $\preAnaRing{A}$ be a pre-analytic ring. An $\underlying{A}$-module $M$ is called \emph{finitely-$\preAnaRing{A}$-free} if there is an isomorphism $M\cong \pfree A S$ for some extremally disconnected set $S$. If more generally $M$ is a direct sum of finitely-$\preAnaRing{A}$-free modules then $M$ is called \emph{$\preAnaRing{A}$-free}.
\end{definition}

\begin{definition}[Presentable modules]
	Let $\preAnaRing{A}$ be a pre-analytic ring. An $\underlying{A}$-module $M$ is called \emph{$\preAnaRing{A}$-presentable} if there is an exact sequence
	$$\bigoplus_{i\in I} \pfree A {S'_i} \to \bigoplus_{j\in J} \pfree A {S_j}\to M\to 0$$
	presenting $M$ in $\mod{\underlying A}$ via $\preAnaRing{A}$-free modules. Write $\pPresentableMod A$ for the (full) subcategory of $\mod{\underlying{A}}$ of $\preAnaRing{A}$-presentable modules.
\end{definition}

As mentioned before, $\preAnaRing{A}$-free modules are not automatically $\preAnaRing{A}$-complete. Solely imposing this condition is however not enough to build a good (derived) theory of $\preAnaRing{A}$-complete modules. Instead, one essentially requires that bound to the right complexes of $\preAnaRing{A}$-free modules are derived $\preAnaRing{A}$-complete.
\begin{definition}[(Derived) completeness]
	Let $\preAnaRing{A}$ be a pre-analytic ring. A complex $C^\bullet \in \derived{\underlying{A}}$ is called \emph{(derived) $\preAnaRing{A}$-complete} if for any extremally disconnected $S$ the morphism
		$$\eRHom_\underlying{A}(\pfree{A}{S}, C^\bullet) \to \eRHom_\underlying{A}(\ufree{A}{S}, C^\bullet)$$
	in $\derived{\Ab}$ is an isomorphism.
\end{definition}

We can observe that derived completeness of an object concentrated in degree $0$ implies ordinary completeness.
\begin{remark}[Derived completeness implies ordinary completeness]
	\label{remark: derived completeness implies ordinary completeness}
	
	Suppose that $X$ is an $\underlying{A}$-module such that $X[0]$ is derived $\preAnaRing{A}$-complete. As $H^0\circ \eRHom_\underlying{A} \cong \eHom_\underlying{A}$ on objects of $\mod{\underlying{A}}$ and since $H^0$ preserves isomorphisms, we obtain that for any $S$ extremally disconnected the induced
		$$\eHom_\underlying{A}(\pfree{A}{S}, X) \to \eHom_\underlying{A}(\ufree{A}{S}, X)$$
	is an isomorphism, hence that $X$ is $\preAnaRing{A}$-complete. It is a priori not clear that the converse holds true! If $\preAnaRing{A}$ is however analytic (see definition \ref{definition: analytic rings}, then it is true a fortiori by theorem \ref{theorem: properties of the category of complete modules}.
\end{remark}

We can now state the definition for an analytic ring. As just mentioned, one needs to impose the condition that bound to the right complexes of $\preAnaRing{A}$-free modules are $\preAnaRing{A}$-complete. Even more, we will require the analogous morphism of condensed abelian groups $\eRHom_\underlying{A}$ to be an isomorphism. This will allow us to control the internal $\intHom_\underlying{A}$'s of $\preAnaRing{A}$-complete modules later on in theorem \ref{theorem: properties of the category of complete modules}.
\begin{definition}[Analytic rings]
	\label{definition: analytic rings}
	
	A pre-analytic ring $\preAnaRing{A}$ is called \emph{analytic} if for any complex
	$$C^\bullet = \left(\dots \to C^{-2} \to C^{-1} \to C^0 \to 0\right)$$
	of $\preAnaRing{A}$-free modules and any extremally disconnected $S$ the induced morphism
	$$\eRHom_\underlying{A}(\preAnaRing A[S], C^\bullet) \to \eRHom_\underlying A(\underlying A[S], C^\bullet)$$
	in $\derived{\condensedAb}$ is an isomorphism.
\end{definition}

\begin{remark}[Free completions are complete]
	If $\preAnaRing{A}$ is analytic and $S$ is extremally disconnected then $\pfree{A}{S}$ is $\preAnaRing{A}$-complete. Indeed, the complex $C^\bullet=\pfree{A}{S}[0]$ is derived $\preAnaRing{A}$-complete and the claim follows by remark \ref{remark: derived completeness implies ordinary completeness}.
\end{remark}

As per usual, we need a good notion of map between analytic rings. Such a map should not(!) simply be a map of the underlying pre-analytic rings.
\begin{definition}[Maps of analytic rings]
	\label{definition: maps of analytic rings}
	
	Let $\preAnaRing A=(\underlying{A}, \pfree{A}{-}, \eta)$ and $\preAnaRing{B}=(\underlying{B}, \pfree{B}{-}, \eta')$ be analytic rings. A \emph{map of analytic rings} $\preAnaRing{A}\to\preAnaRing B$ is a map of the underlying condensed rings $\phi\colon \underlying{A}\to \underlying{B}$ such that for each extremally disconnected $S$ the free module $\pfree B S$ considered as an $\underlying{A}$-module is $\preAnaRing{A}$-complete, \ie the scalar restriction $\scalarRestriction{\phi} \pfree B S$ is $\preAnaRing{A}$-complete.
\end{definition}

The following remark shows how any map of analytic rings induces a map of pre-analytic rings.
\begin{remark}[Maps of analytic rings induce maps of pre-analytic rings]
	\label{remark: maps of analytic rings induce maps of pre-analytic rings}
	
	In the setting of definition \ref{definition: maps of analytic rings}, since $\scalarRestriction{\phi}\pfree B S$ is $\preAnaRing{A}$-complete, the map of condensed sets $\associated S\to \scalarRestriction{\phi}\pfree B S$ factors through $\associated S\to \pfree A S$ via a unique map $\pfree A S \to \scalarRestriction{\phi} \pfree B S$. Hence, a map of analytic rings induces a corresponding map of pre-analytic rings.
\end{remark}

Now to define a map of analytic rings $\preAnaRing{A} \to \preAnaRing{B}$ one has to give a map of underlying rings $\phi\colon \underlying{A}\to\underlying{B}$ and ensure that for any extremally disconnected $S$ the scalar restriction $\scalarRestriction{\phi}\pfree{B}{S}$ is $\preAnaRing{A}$-complete -- this is potentially quite inconvenient. In practice however it is often good enough to specify a map of pre-analytic rings instead.
\begin{proposition}[Maps of analytic rings from maps of pre-analytic rings]
	\label{proposition: maps of analytic rings from maps of pre-analytic rings}
	
	Suppose that $\preAnaRing{A}$ and $\preAnaRing{B}$ are analytic rings and that $(\phi, \epsilon)\colon \preAnaRing{A}\to \preAnaRing{B}$ is a map of pre-analytic rings. If either $\underlying{A}$ is discrete or $\scalarRestriction{\phi}\pfree{B}{\ast}$ is $\preAnaRing{A}$-complete, then $\phi\colon \underlying{A}\to\underlying{B}$ is a map of analytic rings $\preAnaRing{A}\to\preAnaRing{B}$.
\end{proposition}
\begin{shortproof}
	This is weakened version of \cite[Prop. 7.14]{condenseddotpdf}, keeping in mind \cite[Rem. 7.15]{condenseddotpdf}.
\end{shortproof}

We can now state the main result of this section: For any analytic ring $\preAnaRing{A}$ the category of $\preAnaRing{A}$-complete modules is an extremely well-behaved subcategory of $\mod{\underlying{A}}$.
\begin{theorem}[The category of complete $\preAnaRing{A}$-modules]
	\label{theorem: properties of the category of complete modules} 
	
	Let $\preAnaRing{A}$ be an analytic ring.
	\begin{enumerate}
		\item
		The category $\pCompleteMod{A}$ of complete $\preAnaRing{A}$-modules is an abelian category stable under all limits, colimits and extensions in $\mod{\underlying{A}}$. The objects $\pfree A S$ for $S$ extremally disconnected form a (proper) class of compact projective generators of $\pCompleteMod{A}$. The inclusion $\pCompleteMod{A} \to \mod{\underlying{A}}$ admits a left adjoint
			$$\completion{-}{\preAnaRing{A}}\ldef \completionAsTensor{\preAnaRing{A}}{\underlying{A}}{-}\colon\mod{\underlying{A}} \to \pCompleteMod{A}$$
		which is the unique colimit preserving extension of the functor $\ufree A S \mapsto \pfree A S$ defined on the image of $\ufree{A}{-}\colon \extremallyDisconnectedSets\to \mod{\underlying{A}}$ and is called the \emph{$\preAnaRing{A}$-completion} or simply \emph{completion}. If $\underlying{A}$ is commutative, there is a unique symmetric monoidal tensor product $\otimes_\preAnaRing{A}$ on $\pCompleteMod{A}$ for which the completion $\completion{-}{\preAnaRing{A}}$ is (strong) symmetric monoidal. Furthermore, $\pPresentableMod{A}=\pCompleteMod{A}$, that is $\preAnaRing{A}$-presentability and  $\preAnaRing{A}$-completeness agree.
		
		\item
		The functor
			$$\derived{\pCompleteMod{A}} \to \derived{\underlying{A}}$$
		induced by the inclusion $\pCompleteMod{A} \to \mod{\underlying A}$ is fully faithful. Its essential image is stable under all limits and colimits in $\derived{\underlying{A}}$ and is given by those complexes $C^\bullet \in \derived{\underlying{A}}$ for which one of the two equivalent conditions is met:
		\begin{itemize}
			\item
			$C^\bullet$ is derived $\preAnaRing{A}$-complete, that is for all extremally disconnected $S$ the induced map
				$$\eRHom_{\underlying A}(\pfree A S, C^\bullet) \to \eRHom_{\underlying A}(\ufree A S, C^\bullet)$$
			in $\derived{\Ab}$ is an isomorphism. In this case, also the $\eRHom_{\underlying{A}}$s in $\derived{\condensedAb}$ agree.
			
			\item
			For all $n\in\Z$ the $\underlying A$-module $H^n(C^\bullet)$ is $\preAnaRing{A}$-complete.
		\end{itemize}
		Furthermore, the inclusion $\derived{\pCompleteMod{A}} \to \derived{\mod{\underlying{A}}}$ admits a left adjoint
			$$\derivedCompletion{-}{\preAnaRing{A}}\ldef \derivedCompletionAsTensor{\preAnaRing{A}}{\underlying{A}}{-}\colon \derived{\mod{\underlying{A}}}\to\derived{\pCompleteMod{A}}$$
		which is the left-derived functor of $\completion{-}{\preAnaRing{A}} = \completionAsTensor{\preAnaRing{A}}{\underlying{A}}{-}$. If $\underlying{A}$ is commutative, then $\derived{\pCompleteMod{A}}$ admits a unique symmetric monoidal tensor product $\Dotimes_\preAnaRing{A}$ for which $\derivedCompletion{-}{\preAnaRing{A}}$ is (strong) symmetric monoidal.
	\end{enumerate}
\end{theorem}
\begin{proof}
	Most of this follows from lemma \ref{lemma: THE LEMMA} -- once we have proven it in the next section. What remains is that the $\eRHom_\underlying{A}$s in $\derived{\condensedAb}$ agree and that $\derived{\completeMod{\preAnaRing{A}}}$ is symmetric monoidal in case of a commutative $\underlying{A}$.
	
	So suppose that $X^\bullet \in \derived{\completeMod{\preAnaRing{A}}}$.
	Clausen and Scholze argue in the proof of \cite[Prop. 7.5, p.46]{condenseddotpdf} that the claim reduces to the case where $X^\bullet$ is bounded to the right. The author is not sure how they argue this, as he sees no reason why in general $\pfree{A}{S}$ should be $\condensedAb$-enriched compact (which would ensure the commutation with the homotopy colimit $\colimit_{n\in \N} \cTruncation{\le n}X^\bullet$) or why $\eRHom_\underlying{A}$ should commute with limits (which would ensure the commutation with the limit $\limit_{n\in \N} \sTruncation{\le n}X^\bullet$) -- only commutation with homotopy limits seems ensured. We thus have to make it an \assumptionPreWord that we can reduce to the case that $X^\bullet$ is bounded to the right. As $\eRHom_A$ is compatible with shifts we can further assume that $X^\bullet$ is concentrated in non-positive degrees. Then $X^\bullet$ admits a left resolution by a complex $C^\bullet$ of $\preAnaRing{A}$-free modules that is concentrated in non-positive degrees too. By definition of analyticity of $\preAnaRing{A}$ we then have
		$$\eRHom_{\underlying{A}}(\pfree{A}{S}, X^\bullet) \cong \eRHom_\underlying{A}(\pfree{A}{S}, C^\bullet) \isorightarrow \eRHom_\underlying{A}(\ufree{A}{S}, C^\bullet) \cong \eRHom_{\underlying{A}}(\ufree{A}{S}, X^\bullet)$$
	in $\derived{\condensedAb}$ which shows the claim.
	
	It remains to argue the existence of the monoidal structure on $\derived{\completeMod{\preAnaRing{A}}}$ for which the derived completion is (strong) symmetric monoidal in case that $\underlying{A}$ is commutative. For this we will simply reference the proof given by Clausen and Scholze for \cite[Prop. 7.5]{condenseddotpdf}.
\end{proof}

\begin{notation}
	If $\preAnaRing{A}$ is an analytic ring we will abbreviate
		$$\mod{\preAnaRing{A}}\ldef \pCompleteMod{\preAnaRing{A}}\text{ and }\derived{\preAnaRing{A}}\ldef \derived{\pCompleteMod{A}}$$
	and refer to these categories as the \emph{category of $\preAnaRing{A}$-modules} and the \emph{derived category of $\preAnaRing{A}$-modules} respectively.
\end{notation}

\begin{remark}[$\Dotimes_\preAnaRing{A}$ and $\otimes_\preAnaRing{A}$]
	\label{remark: when is the tensor product the derived tensor product?}
	
	There is a notational subtlety: For a general analytic ring $\preAnaRing{A}$ the functor $\Dotimes_\preAnaRing{A}$ might not be the left derived functor of $\otimes_\preAnaRing{A}$!
	
	Clausen and Scholze state in \cite[Warning 7.6]{condenseddotpdf} that this however holds true in any instance know to them. Furthermore, they state that one might equivalently check that for any two extremally disconnected sets $S$ and $T$ the complex of $\underlying{A}$-modules $\pfree{A}{S\times T}\ldef \nerve{\pfree{A}{S_\bullet}}$ (obtained by simplicially resolving the non-extremally disconnected set $S\times T$ by extremally disconnecteds, \ie choosing a simplicial hypercover $S_\bullet$ of $S\times T$ of extremally disconnected sets $S_n$, giving a simplicial $\underlying{A}$-module $\pfree{A}{S_\bullet}$ and attaching to it by the Dold--Kan correspondence a connective complex $\nerve{\pfree{A}{S_\bullet}}$) is concentrated in degree $0$.
	
	We will see in remark \ref{remark: free completions for profinite sets} that we can -- for all analytic rings introduced in this thesis -- safely ignore the issue and always assume $\Dotimes_\preAnaRing{A}$ to be the left derived functor of $\otimes_\preAnaRing{A}$.
\end{remark}

Now knowing, that complete modules give rise to extremally well-behaved abelian categories, we should also investigate how one can transport complete modules along a map of analytic rings.
\begin{proposition}[Scalar restriction and extension of complete modules]
	\label{proposition: scalar restriction and extension of complete modules}
	
	Let $\phi\colon \preAnaRing{A}\to \preAnaRing{B}$ be a map of analytic rings.
	\begin{enumerate}
		\item
		The scalar restriction of any $\preAnaRing{B}$-complete module is $\preAnaRing{A}$-complete and the induced functor $\scalarRestriction{\phi}\colon \mod{\preAnaRing{B}}\to\mod{\preAnaRing{A}}$ admits a left adjoint
			$$\completedScalarExtension{\phi}{\preAnaRing{B}}\ldef \preAnaRing{B}\otimes_{\preAnaRing{A}}-\colon \mod{\preAnaRing{A}}\to\mod{\preAnaRing{B}}$$
		that is the unique colimit preserving extension of $\pfree{A}{S}\mapsto \pfree{B}{S}$ and is called the \emph{completed scalar extension}.
		
%		Scalar restriction $\scalarRestriction{\phi}\colon \mod{\underlying{B}}\to\mod{\underlying{A}}$ restricts to a functor $\mod{\preAnaRing{B}}\to\mod{\preAnaRing{A}}$. Scalar extension $\scalarExtension{\phi}\colon \mod{\underlying{A}}\to\mod{\underlying{B}}$ gives rise to a functor
%			$$\mod{\underlying{A}} \longxrightarrow{\scalarExtension{\phi}} \mod{\underlying{B}} \longxrightarrow{\preAnaRing{B}\otimes_{\underlying B} -} \mod{\preAnaRing{B}}$$
%		which factors over the $\preAnaRing{A}$-completion $\mod{\underlying{A}}\to\mod{\preAnaRing{A}}$ via a functor denoted 
%			$$\completedScalarExtension{\phi}{\preAnaRing{B}}\ldef\preAnaRing{B}\otimes_\preAnaRing{A} -\colon \mod{\preAnaRing{A}}\to\mod{\preAnaRing{B}},$$
%		that is the unique colimit preserving extension of $\pfree{A}{S}\mapsto \pfree{B}{S}$ and is called the \emph{scalar extension}. Furthermore $\completedScalarExtension{\phi}{\preAnaRing{B}}$ is left adjoint to $\scalarRestriction{\phi}\colon \mod{\preAnaRing{B}}\to\mod{\preAnaRing{A}}$.
		
		\item
		The scalar restriction of any complex in $\derived{\preAnaRing{B}}$ is $\preAnaRing{A}$-complete and the induced functor $\scalarRestriction{\phi}\colon \derived{\preAnaRing{B}}\to\derived{\preAnaRing{A}}$ admits a left adjoint
			$$\derivedCompletedScalarExtension{\phi}{\preAnaRing{B}}\ldef \preAnaRing{B}\Dotimes_\preAnaRing{A} -\colon \derived{\preAnaRing{A}}\to\derived{\preAnaRing{B}}$$
		that is the left derived functor of the completed scalar extension $\completedScalarExtension{\phi}{\preAnaRing{B}}$.
		
%		The composite functor
%			$$\derived{\underlying{A}} \longxrightarrow{\LeftD\scalarExtension{\phi}} \derived{\underlying{B}} \longxrightarrow{\preAnaRing{B}\Dotimes_{\underlying B} -} \derived{\preAnaRing{B}}$$
%		factors over $\derived{\preAnaRing{A}}$ via the left derived scalar extension, denoted
%			$$\preAnaRing{B}\Dotimes_\preAnaRing{A} -\colon \derived{\preAnaRing{A}}\to \derived{\preAnaRing{B}}.$$
	\end{enumerate}
\end{proposition}
\begin{proof}
	Denote by $i_\preAnaRing{A}\colon \mod{\preAnaRing{A}}\to\mod{\underlying{A}}$ and $i_\preAnaRing{B}\colon \mod{\preAnaRing{B}}\to\mod{\underlying{B}}$ the exact inclusions.
	\begin{enumerate}
		\item
		By definition of maps of analytic rings, for any extremally disconnected $S$, the scalar restriction $\scalarRestriction{\phi}\pfree{B}{S}$ is $\preAnaRing{A}$-complete. As scalar restriction commutes with arbitrary colimits and $\mod{\preAnaRing{B}}$ is generated, under colimits, by such $\pfree{B}{S}$, we obtain that scalar restriction restricts to a functor $\mod{\preAnaRing{B}}\to\mod{\preAnaRing{A}}$. Hence, we obtain the commutative square
		$$\begin{tikzcd}
			\mod{\preAnaRing{B}} \ar[r, "\scalarRestriction{\phi}", swap]\ar[d, "i_\preAnaRing{B}"] &\mod{\preAnaRing{A}}\ar[d, "i_\preAnaRing{A}", swap]\\
			\mod{\underlying{B}}\ar[r, "\scalarRestriction{\phi}"]&\mod{\underlying{A}}.
		\end{tikzcd}$$
		Observe that the inclusions $i_\preAnaRing{A}$ and $i_\preAnaRing{B}$ as well as the scalar restriction $\scalarRestriction{\phi}\colon \mod{\underlying{B}}\to \mod{\underlying{A}}$ admit left adjoints given by the completions $\completion{-}{\preAnaRing{A}}$ and $\completion{-}{\preAnaRing{B}}$ as well as scalar extension $\scalarExtension{\phi}\colon \mod{\underlying{A}}\to\mod{\underlying{B}}$. Define the functor
			$$\completedScalarExtension{\phi}{\preAnaRing{B}}\colon \mod{\preAnaRing{A}}\to\mod{\preAnaRing{B}}$$
		as the composition $M\mapsto \completion{\scalarExtension{\phi} i_\preAnaRing{A}M}{\preAnaRing{B}}$. Observe, that for any $\preAnaRing{A}$-complete module $M$ and any $\preAnaRing{B}$-complete module $N$ we have a natural isomorphism
			$$\Hom(\completedScalarExtension{\phi}{\preAnaRing{B}}M, N) \cong \Hom(i_\preAnaRing{A}M, \scalarRestriction{\phi}i_\preAnaRing{B} N) = \Hom(i_\preAnaRing{A}M, i_\preAnaRing{A}\scalarRestriction{\phi} N) \cong \Hom(M, \scalarRestriction{\phi} N)$$
		by adjointness, by commutativity of the above square and by full faithfulness of $i_\preAnaRing{A}$. In particular, $\completedScalarExtension{\phi}{\preAnaRing{B}}$ is the unique (up to unique isomorphism) left adjoint of $\scalarRestriction{\phi}\colon \mod{\preAnaRing{B}}\to\mod{\preAnaRing{A}}$. Thus, all functors in the above commutative square admit a left adjoint and the induced diagram
		$$\begin{tikzcd}
			\mod{\preAnaRing{B}} &\mod{\preAnaRing{A}}\ar[l, "\completedScalarExtension{\phi}{\preAnaRing{B}}", swap]\\
			\mod{\underlying{B}}\ar[ur, Rightarrow, "\sim" {sloped, near start}]\ar[u, "{\completion{-}{\preAnaRing{B}} = \preAnaRing{B}\otimes_\underlying{B} -}"] &\mod{\underlying{A}}\ar[l, "\scalarExtension{\phi}"]\ar[u, "{\completion{-}{\preAnaRing{A}} = \preAnaRing{A}\otimes_\underlying{A} -}", swap].
		\end{tikzcd}$$
		of left adjoints must commute up to unique isomorphism as well, as the composition of two left adjoints is a left adjoint of the composition of the original functors. In particular,
			$$\completedScalarExtension{\phi}{\preAnaRing{B}}\pfree{A}{S} \cong  \completedScalarExtension{\phi}{\preAnaRing{B}}\completion{\ufree{A}{S}}{\preAnaRing{A}} \cong \completion{\scalarExtension{\phi}\ufree{A}{S}}{\preAnaRing{B}} = \completionAsTensor{\preAnaRing{B}}{\underlying{B}}{\underlying{B}\otimes_\underlying{A}\ufree{A}{S}} \cong \completionAsTensor{\preAnaRing{B}}{\underlying{B}}{\ufree{B}{S}} \cong \pfree{B}{S}$$
		so $\completedScalarExtension{\phi}{\preAnaRing{B}}$ is the unique colimit preserving extension of $\pfree{A}{S} \mapsto \pfree{B}{S}$ as claimed.
		
		\item 
%		Observe that $i_\preAnaRing{A}$ and $i_\preAnaRing{B}$ are exact, since $\mod{\preAnaRing{A}}$ and $\mod{\preAnaRing{B}}$ are abelian subcategories of $\mod{\underlying{A}}$ and $\mod{\underlying{B}}$ respectively.
		Scalar restriction $\scalarRestriction{\phi}$ is exact. Recall that for an analytic ring, a complex in the derived category is complete if and only if all of its cohomology groups are complete. Thus, if $C^\bullet\in\derived{\preAnaRing{B}}$ then for each $n\in \Z$ the $\underlying{B}$-module $H^n(C^\bullet)$ is $\preAnaRing{B}$-complete. But as $\scalarRestriction{\phi}$ is exact we obtain that $H^n(\scalarRestriction{\phi}C^\bullet) \cong \scalarRestriction{\phi}H^n(C^\bullet)$ is $\preAnaRing{A}$-complete by the first part and hence $\scalarRestriction{\phi}C^\bullet \in \derived{\preAnaRing{A}}$, that is $\scalarRestriction{\phi}$ restricts to a functor $\derived{\preAnaRing{B}}\to\derived{\preAnaRing{A}}$. The remaining proof is then analogous to the first part.
	\end{enumerate}
\end{proof}

\begin{remark}[Notational compatibility]
	Recall that any condensed ring $A$ can be considered as a pre-analytic ring, denoted $A$ as well. This pre-analytic ring $A$ is clearly analytic and any module is automatically complete. Then the notation $\mod{A}$ for complete modules, as well as the notation $\preAnaRing{A}\otimes_{A} -$ for any analytic ring $A\to \preAnaRing{A}$ over $A$ and $A\otimes_\preAnaRing{A}-$ for any analytic ring $\preAnaRing{A}\to A$ under $A$ are compatible with this identification.
\end{remark}

We can further observe the following.
\begin{lemma}[Completed scalar extension is monoidal]
	\label{lemma: completed scalar extension is monoidal}
	
	Let $\phi\colon \preAnaRing{A}\to \preAnaRing{B}$ be a map of analytic rings. Then completed scalar extension
		$$\completedScalarExtension{\phi}{\preAnaRing{B}}\colon \mod{\preAnaRing{A}}\to\mod{\preAnaRing{B}}$$
	and derived completed scalar extension
		$$\derivedCompletedScalarExtension{\phi}{\preAnaRing{B}}\colon \derived{\preAnaRing{A}}\to\derived{\preAnaRing{B}}$$
	are (strong) symmetric monoidal.
\end{lemma}
\begin{proof}
	We will prove the claim only for $\completedScalarExtension{\phi}{\preAnaRing{B}}$ as the other case is analogous. Observe first that for any two extremally disconnected $S$ and $T$ we have
		$$\pfree{A}{S}\otimes_\preAnaRing{A} \pfree{A}{T} \cong \ufree{A}{S}^\completionSymbol{\preAnaRing{A}} \otimes_\preAnaRing{A} \ufree{A}{T}^\completionSymbol{\preAnaRing{A}}\cong\completion{\ufree{A}{S}\otimes_\underlying{A}\ufree{A}{T}}{\preAnaRing{A}} \cong \ufree{A}{S\times T}^\completionSymbol{\preAnaRing{A}}$$
	as $\completion{-}{\preAnaRing{A}}$ is (strong) monoidal. Similarly, $\pfree{B}{S}\otimes_\preAnaRing{B} \pfree{B}{T} \cong \ufree{B}{S\times T}^\completionSymbol{\preAnaRing{B}}$ for $\preAnaRing{B}$. Recall further, from the proof of proposition \ref{proposition: scalar restriction and extension of complete modules}, that $\completedScalarExtension{\phi}{\preAnaRing{B}}\circ \completion{-}{\preAnaRing{A}} \cong \completion{-}{\preAnaRing{B}}\circ \scalarExtension{\phi}$. We now obtain
	\begin{align*}
		\completedScalarExtension{\phi}{\preAnaRing{B}}\mleft(\pfree{A}{S}\otimes_\preAnaRing{A}\pfree{A}{T}\mright)
		&\cong \completedScalarExtension{\phi}{\preAnaRing{B}} \mleft(\ufree{A}{S\times T}^\completionSymbol{\preAnaRing{A}}\mright)\\
		&\cong \mleft(\scalarExtension{\phi}\ufree{A}{S\times T}\mright)^\completionSymbol{\preAnaRing{B}}\\
		&\cong \ufree{B}{S\times T}^\completionSymbol{\preAnaRing{B}}\\
		&\cong\pfree{B}{S}\otimes_\preAnaRing{B}\pfree{B}{T} \cong \completedScalarExtension{\phi}{\preAnaRing{B}}\pfree{A}{S} \otimes_\preAnaRing{B} \completedScalarExtension{\phi}{\preAnaRing{B}}\pfree{A}{T}
	\end{align*}
	which proves that $\completedScalarExtension{\phi}{\preAnaRing{B}}$ is (strong) monoidal on generators. As $\completedScalarExtension{\phi}{\preAnaRing{B}}$ as well as $\otimes_\preAnaRing{A}$ and $\otimes_\preAnaRing{B}$ commute with arbitrary colimits, this already shows the claim.
\end{proof}

